\textbf{Date:} November 18, 2026
\\\textbf{Interviewer:} Lucas DeBoer
\\\textbf{Participant:} P4, an individual who has been blind since birth (near total sight loss, can see silhouettes)
\\\textbf{Format:} In-person conversation, about 30 minutes long
\\\textbf{Editing:} The entire conversation was transcribed and recorded using Apple's Voice Memo app. Light editing was done afterwards to fix transcription mistakes and keep the transcript readable by keeping Lucas' one-word replies (``Yeah,'' ``Yep,'' ``Gotcha'') inside P4's turns as \textit{[Lucas: Yeah.]}. Personally identifying information and information not relevant to the interview were also removed or redacted. When P4's phone spoke (for VoiceOver), the transcribed text is marked by \textit{[phone speech]} or removed if confidential.
\\\textbf{Recruitment:} A family friend of the interviewer, recruited directly
\\\textbf{Consent:} P4 verbally consented to the interview and its recording prior to the recording beginning.
\\\textbf{Anonymization:} P4 is not named. All identifying details have been removed.


\bigskip
\hrule
\bigskip

\textbf{Lucas:} Yeah. Uh, okay, I gave you the introduction as to what we're doing. So we're just looking for what works well for you, what sort of things that are frustrating for you. Obviously, this isn't a test, there's no right or wrong answers, just people looking to see what you feel. And you can skip any questions you want. Okay. So if you could tell me a little bit about yourself and how technology fits into your day, as much as you're comfortable, how does your disability affect how you use a computer and has that changed over time?

\textbf{P4:} It's changed a lot over time. I'll sort of give it a little background about myself so everybody kind of knew. We live on a chicken farm here, and I graduated from the W.\ Ross McDonald School, which is the School for the Blind in Ontario. And then I went to George Brown College and took a piano tech course. That was all about tuning and building or repairing pianos. And so, uh, computer technology has changed a lot in all of that time. When I was in the early grades, in high school, we got talking calculators, and that was a big deal. But they were, the voice was very monotone. Um, and it could speak the numbers and the signs. But that's it, though. And then, uh, oh, and then we got Commodore 64s. And but they didn't have a speech program, so they were kind of useless at the school, except for the print users. But then we got the Apple 2 E's, and they had a voice program. They had a screen reader program, which was a voice thing. So up until then, there was the Versa Braille and the Braille link, and they had a Braille display. But this was really the first computer that spoke. It had a screen reader program. So, um, and now then, uh, it's updated further from there. And Microsoft had a version called Jaws, which was pretty basic in the beginning, but it's a lot more, you can do lots with it now. And then Apple had something that they called VoiceOver, you know? And so, um, I use the Apple only because it seems to be, um, well, for one thing, you can vibe with stuff for less money. Because the disability access is included in every product that they have, and it's the same. So you can go from an iPad to a phone or to a Mac and the commands are very similar. \textit{[Lucas: Yeah.]} So it's just more convenient. And I think, um, for what I do, it actually works out a little better. And so, um, the only thing that I'm really running with the computer here is the control in the chicken barn. \textit{[Lucas: Yeah.]} And so, um, the phone reads what's on the screen there, and then, um, I go to the item that I want. I select what I want, and then I can edit, and then I press return, and it locks in. \textit{[Lucas: Yeah, it does its thing, as you would know.]} Yep. And, um, so, uh, Maximus is the controller that I use in it. It's the only one that will work with my phone, and it happens to be an Apple-based product.

\textbf{Lucas:} Gotcha. So it works well.

\textbf{P4:} \textit{[P4's Chicken Salesman]} wasn't too happy.

\textbf{Lucas:} I was going to say, because you had talked how you used the Genius before.

\textbf{P4:} Yeah, but it doesn't. And it doesn't\ldots{} Well, it sort of kind of, it would be a lot of work, I'm told, to get it to work. And so it's just too, this works pretty good. So the other thing that I use the computer for is a little bit of word processing. I do email, and I do texting, and so I depend, when I'm typing an email, I do depend on the screen reader program to let me know what I've written, if I want to, and so there's different settings for that. You can have it speak every letter as you go, or you can just have it speak the words, or you can have it just off and then turn it on and read the thing when you're done. \textit{[Lucas: Yeah.]} You know, if you want to proofread your email before you send it, like, that's not necessary, you know, which happens to me once in a while. Depends where it's going. \textit{[Lucas: Yep.]} So, um, my brother, he's using technology all the time, and he would be probably a better interview about this than me, actually. Um, he's quite up on what's out there. He works for something called \textit{[P4 brother's palce of work]}, and he's been there a long time, and he's done well there, and he does some secretarial sort of stuff. He's answering the phone there, and he's also taking notes at meetings, that kind of thing. And, I don't know what else he does there, but they keep him. They haven't fired him.

\textbf{Lucas:} Fair enough. Um, so you answered, like, three of my questions in a row which is perfect. So you're typically using the computer for word processing, emails, texting, and controlling the chicken barn.

\textbf{P4:} Correct. Yeah, it's mostly that.

\textbf{Lucas:} Um, you mentioned you use VoiceOver or screen reader. Is there any other accessibility tools, features, or voice assistance that you use?

\textbf{P4:} No, that's it.

\textbf{Lucas:} Now, I haven't actually tried to use it myself. Is VoiceOver different from a screen reader? Is it a screen reader?

\textbf{P4:} Well, so VoiceOver is the name of the screen reader program. Just like the Microsoft version is Jaws. \textit{[Lucas: Yeah.]} And there's a couple other ones out there. But the Apple one is called VoiceOver.

\textbf{Lucas:} And so VoiceOver is different than Siri.

\textbf{P4:} Yeah, of course. So, um, VoiceOver is just the screen reader program, and then I have to, um, I have to tap the screen, uh, do a double tap if I want to select something. Okay? Or, which I can show you.

\textbf{Lucas:} Sure, I'd like to see how this all works.

\textbf{P4:} So there, I just got a text from \textit{[P4's son, friend of interviewer]}, you know him?

\textbf{Lucas:} I heard of him once or twice.

\textbf{P4:} Hey, so, yeah\ldots{} We'll just go to\ldots{} Well, go to\ldots{} So, it says\ldots{} So, I touch the phone icon, so it says phone.

\textit{[phone speech: ``Messages. 3 unread messages.'']}

\textbf{P4:} Okay, so I've got three unread messages that tells me. \textit{[Lucas: Yep.]} And then, so I have to double tap, okay?

\textit{[phone speech: Starts reading a previous message]}

\textbf{P4:} Oh, so he's talking about a grain cart or something. Today he says what he delivered. So if I want to go back, okay, then I need to swipe like this.

\textit{[phone speech: Read back other messages]}

\textbf{Lucas:} The swiping back and forth, that's controlled by VoiceOver? That's not, like, a separate program on top of it?

\textbf{P4:} Right, that's part of VoiceOver, and that's how, so now this is kind of a long text. There's a lot of information went by pretty quick. \textit{[Lucas: Yeah.]} So I can speed that up and slow it down if I want to. That's all part of VoiceOver. And then, so now, I can also do, um, do this, so I'm gonna do, like, a rotor thing. \textit{[Lucas: Yeah.]} So that'll speak the characters now, so I'm going the other way.\\
\textit{[phone speech: ``Action. Activate. Words.'']}

\textbf{P4:} There's words. Okay, so now we can do line at a time.

\textit{[phone speech: reading the listing line by line]}

\textbf{P4:} Right, so\ldots{} Now, so 30.35 or 32 or whatever. So I didn't quite get all that. \textit{[Lucas: Yeah]} Um, so I can go to characters, okay?

\textit{[phone speech: reading a line of listing character by character]}

\textbf{P4:} Okay, so that's the\ldots{} So I can get that out of there. Okay, so I'm done with that. So now I want to go back to lines, so I just twist it back. So that's the end of that. So now, okay, so I'm out of that text, so now I want to go to the next one. \textit{[Lucas: Yeah.]} So I just swipe.

\textit{[phone speech: reading other messages]}

\textbf{P4:} So then, if I wanted to answer him, or I did\ldots{} But then I get the message there on the bottom. And then you tap it twice. And then there's the keyboard. \textit{[Lucas: Yep.]} And then now I can dictate to it if I want to. \textit{[Lucas: Yep.]} So I'll do that for a minute.\\
\textit{[dictating]}\\
Um, and if I speak the punctuation, sort of, um, a little bit lesser than the text, then it does do the punctuation. Okay, but you gotta be\ldots{} Sometimes it might type out ``period.'' \textit{[Lucas: Yeah, yeah.]} Or instead of ``comma,'' I might say ``come on,'' you know, some sort of something. Okay. So now if I want to send this to him, we'll just finish this off.\\
\textit{[dictating]}.\\
Send. So, it says send.

\textbf{Lucas:} Okay, so you still had to, you kind of have to click where the message is there. Like at the beginning, you were clicking on the bottom. If you had clicked on the bottom, you had to move your way all the way up to the top, if you didn't know to click at the top, say.

\textbf{P4:} Oh, well, yeah, that's the one problem, um, so charts and graphs and that kind of thing, spreadsheets. \textit{[Lucas: Yeah.]} Um, are a pain. I guess. And so this thing will tell me, you know, row 5, column 6. \textit{[Lucas: Yeah, yeah.]} Or something. And so, um, that kind of thing. Yeah, it's just, it's a tough go. That's what I've got \textit{[P4's wife]} for. Like people say, ``Hey, Siri,'' I say, ``Hey, \textit{[P4's wife]}.''

\textbf{Lucas:} Gotcha.

\textbf{Lucas:} Okay, so you mentioned that VoiceOver is different than Siri. Do you use Siri much or no? Is Siri kind of useless?

\textbf{P4:} Uh, well, yeah, I guess I don't really. So what I do sometimes, I'll just hold down the, the whole button. And then so if I'm looking for music, then I'll just say, ``YouTube, [song title].'' So I'd hold down the whole button and say, ``YouTube, You Are the Wind Beneath My Wings, Bette Midler.'' So then that pulls that up. Okay, right. So, like, I never really say, ``Hey, Siri.'' \textit{[Lucas: Yeah, yeah, yeah.]} You know, and all of that. Um, like people use Alexa to play the radio station. I just\ldots{} I'll just say, no, I never say ``Hey, Siri.'' So if I want to listen to 680 News, I'll just open Siri and say ``Play 680 News Toronto.'' and it pulls it up. Okay, and then, oh, sometimes I have to say, use Apple Music.

\textbf{Lucas:} Oh, okay. Um, so I'm kind of focused on VoiceOver, 'cause that seems to be what you use the most. How does VoiceOver work for you?

\textbf{P4:} Like, overall, pretty good.

\textbf{Lucas:} You mentioned that spreadsheets and charts and whatnot are a bit of a pain.

\textbf{P4:} Yeah, so overall, it works pretty good.

\textbf{Lucas:} Yeah, yeah. So if you need help with computer tasks, such as completing an action, setting up a tool, what do you usually do?

\textbf{P4:} Well, I usually ask \textit{[P4's wife]}. \textit{[Lucas: Okay.]} See, and then try again.

\textbf{Lucas:} So say Susan wasn't around, There's no way for me to get around that. You're just at a lock.

\textbf{P4:} Correct, okay, so I can show you a problem. Actually, there's also in it, just kind of clean the screen.

\textbf{Lucas:} That's perfect. You're, like, guessing all of my questions. This is awesome.

\textbf{P4:} I know. I'm a very valuable research subject.

\textbf{Lucas:} You know it.

\textbf{P4:} So we're gonna go to the Maximus app here. Yeah? So it wants to do an update there, okay? And for some reason, that's not responding to me. Okay, so, like, if I do this \textit{[tries to user VoiceOver]}. It tells me that.

\textit{[phone speech: reads the update message on the screen]}

\textbf{P4:} So then, right now, it hasn't given me any instructions as to what to do. Okay, so now we're on the bottom there. So, it hasn't actyually told me how to perform the update.

\textbf{Lucas:} Gotcha. Yeah, 'cause they're assuming you can just see that and click whatever they want.

\textbf{P4:} Right. So for some reason, and maybe the update will fix that. I don't know. But right now, um, you know, it's not, so that'll be one of the things I say, to either that or call the guys that I bought the thing from.

\textbf{Lucas:} Okay.

\textbf{P4:} And they're pretty good with that sort of thing.

\textbf{Lucas:} But obviously that's a step away that takes time and whatnot.

\textbf{P4:} Well, it's right. It's annoying and it's nice to, so this showed up, um, a couple of nights ago, and I've been a bit busy the last couple days. \textit{[Lucas: Yeah.]} And haven't been able to do it, but I can, I can bypass it as you can see there. \textit{[Lucas: Yeah.]} But, um, so that's a problem.

\textbf{Lucas:} Yep. Gotcha. When you're doing a task, how does your computer let you know what's happening, whether it understood you, what's on the screen, whether something's working? Um, so it's just reading it back to you. There's no sort of other way that\ldots{}

\textbf{P4:} Right, yeah, it's just reading back. It'll read the progress bar, um, and update that way if I want to, um, I can access that.

\textbf{Lucas:} Yeah. And it is just reading exactly what's on the screen. It's not kind of giving a summary based off of context.

\textbf{P4:} No, it's just reading what's on the screen. So if there's a progress bar that says, you know, 50\% or 60\% kind of thing, it'll just read that.

\textbf{Lucas:} Okay. Um, thinking about the last time you did something online that mattered to you, or was important, just sending an email, submitting a form, making a purchase. How did you make sure that it was right before it was finished? Or how can you make sure that it's right for you, that it's finished?

\textbf{P4:} Oh, um\ldots{} So I do some online banking. \textit{[Lucas: Yeah.]} And then, um, before I send an e-Transfer, trying to think of what it all says on the screen, but there is a point where it says details, okay? And then, um, it'll give me the email address that it's going to and that kind of thing. So then it'll say that the fields are there. And as I do them. But there's another problem there, though, I'm sort of thinking about it, I'm telling you that.

\textbf{Lucas:} Perfect.

\textbf{P4:} So, what RBC does, oh, and CFO \textit{[Chicken Farmers of Ontario, refering to their website]}, it's not accessible. I actually got someone whose job it is to figure out these things. \textit{[Lucas: Oh yeah.]} So I took my computer to him and said that this is what I need to do. \textit{[Lucas: Yeah.]} And he looked at it for like two minutes and said, quote, ``You don't have a hope in hell.'' So then, of course, the CFO, they sent somebody out to my house for the day to figure this out, to teach me how to do this. \textit{[Lucas: Yeah.]} I just think that that was, and it was kind of hopeless.

\textbf{Lucas:} Gotcha. So you mentioned there's a problem with the online banking.

\textbf{P4:} So the problem with the online banking that RBC does not want to do anything about, um, is that if I want to send an e-Transfer, um, if the e-Transfer is too close to the maximum amount, and this is only in the personal account, not in the business account, and don't ask me why, then it'll give me a screen that wants my PIN number. \textit{[Lucas: Yeah.]} But that screen doesn't respond to the number pad that is ordinarily there with the banking thing. Okay, so, um, and I'm not sure why. \textit{[Lucas: Yeah.]} It doesn't, because the rest of it works. Well, so there are, as you probably know, there's codes, the standards, this stuff is supposed to work.

\textbf{Lucas:} Yes, yup.

\textbf{P4:} For the disabled. And so, the RBC banking app does work, for the most part, except for, like, a little glitch like that once in a while. Then, like, you know, you're trying to do your thing, and then all of a sudden you get that screen that's got the number pad on it, but it looks a bit different than the other one. \textit{[Lucas: Gotcha.]} And then all of a sudden I just get in a loop that I can't get out of. And so, like, stuff like that is, and the frustrating part is when you tell them, of course, they don't understand. And then you say, ``Oh, yeah, okay, we'll look into it.'' \ldots{} Yeah, well how you doing with that?

\textbf{Lucas:} But VoiceOver in general does work with banking? Because when I had talked with someone in another interview, the voice command that he uses doesn't work with online banking.

\textbf{P4:} Oh, okay, so what was that? Samsung then or something?

\textbf{Lucas:} Yes, it was.

\textbf{P4:} Oh, yeah. See, Apple is better than Samsung that way. The screen reader program is better. And there's where, like, \textit{[P4's brother]} kind of knows what's really going on.

\textbf{Lucas:} So you mentioned websites or apps, sometimes they're not coded up to the standard that they should be.

\textbf{P4:} Yeah, like the CFO website, you know, it's not coded to\ldots{} Yeah, it doesn't work. Actually, that said, they just did an update, so I need to probably check that. \textit{[Lucas: Yeah.]} But up until now, it's been quite hopeless. \textit{[Lucas: Yeah.]} And the other thing that, so that, that, um, well, there's lots wrong with the CFO. But, um, so the other thing that I tried to get the CFO to do was to put the information that we, you know, might need, such as the news component, to just put it at the top somewhere. Don't bury it too far in. \textit{[Lucas: Yes.]} Because that's when it gets to be a problem. Right? And then so, um, if I go to something like the Real Ag newsletter, oh, yeah, okay, so there's another problem. If I go to the Real Ag newsletter, I click on this story with my phone. Um, and for some reason, um, I can't find the, uh, the play video button. And so, like, lots of times they, uh, on the Real Ag newsletter, there's a summary of an interview and then they post a video of it as well. And you can listen to the video and, you know, a lot of the time the summary's just good enough. But sometimes, though, you want to listen to the whole thing. \textit{[Lucas: Yeah.]} And so if I could say, ``Play the video that corresponds with this,'' \textit{[Lucas: Yeah, yeah, yeah.]} you know, that'd be helpful, kind of thing.

\textbf{Lucas:} Perfect. And now, does VoiceOver mess up? Or since it's not really deciding anything on its own, it just does what you say.

\textbf{P4:} Yeah, it just does what I say. \textit{[Lucas: Okay.]} Or, well, it just does. It's all about like the sliding and then the double tapping, kind of thing. So, um, yeah, it doesn't really mess up. The only thing, I guess, that it messes up in is that it decides it doesn't know what to do with it. So, like, charts and graphs are a problem. Or the other thing that is a problem is if somebody sends me a scanned document then the computer won't read it.

\textbf{Lucas:} Oh, okay. Interesting. Now, your phone and your computer, your Mac, you kind of mentioned that earlier, they're fairly connected, like fairly interchangeable.

\textbf{P4:} Yeah.

\textbf{Lucas:} Um, is there anything that you want to do on a computer that you currently avoid or have someone else do?

\textbf{P4:} Oh, um, yeah, some of the\ldots{} I've never learned QuickBooks. So I get \textit{[Related family member]} to do it.

\textbf{Lucas:} Okay. Is there anything else that you would like to share that you think's important?

\textbf{P4:} And sort of, like, the CFO data that I need to be entering, like when it's shipping time and stuff like that. Like, it would be better if I could do it on my own. So, but because of the, you know, CFO, their programming is just quite hopeless.

\textbf{Lucas:} Yeah. So really it comes down to more that the websites and such that you want to use are not created properly.

\textbf{P4:} Yeah, that's right.

\textbf{Lucas:} Because VoiceOver itself is fairly competent. It's more just the stuff that you'd be using VoiceOver on isn't done well.

\textbf{P4:} Correct.

\textbf{Lucas:} Gotcha.

\textbf{P4:} And so even though, like, there are tests for them to run it through. \textit{[Lucas: Yeah.]} But a lot of the time, those are not accurate real life tests. So somebody may run it through there and it may be reading something. But that doesn't mean it's reading what you needed to read.

\textbf{Lucas:} Ah, yeah. Kind of the context behind it.

\textbf{P4:} Yeah, or, you know, that's a thing.

\textbf{Lucas:} Yeah. Um\ldots{} Is there anything else you want to share regarding your experience? Like, we've focused a lot on bad stuff. Is there anything that VoiceOver or other stuff does really well?

\textbf{P4:} Um, well, yeah, like the way I can use the Maximus, um, for the most part. I really like that. And the fact that I can text and all of that, an email. \textit{[Lucas: Yeah.]} Like, it means that I can do committee work and all of that kind of thing. So, I appreciate that.

\textbf{Lucas:} Great, now sometime next year we will probably need people to test out what we have built, are you open to being a tester?

\textbf{P4:} Oh, yes. Definitely.
